% Context from chapters/wavelets.tex; for mathematical reference, not compilation.
.
	}
	If there exist $0 < a \leq b <+\infty$ such that $a \leq A(\om) \leq b$ almost everywhere, then $\phi$ defined by
	\eq{
		\hat\phi(\om) = \frac{\hat\th(\om)}{\sqrt{A(\om)}}
	}
	is such that $\{\phi(\cdot-n)\}_{n \in \ZZ}$ is an orthonormal basis of $\overline{\Span}\{\th(\cdot-n)\}_{n\in\ZZ}$.
\end{prop}


\begin{figure}[tbp]
\centering
\BookDrawing{tikz/wavelets-spline-orthogonalization}
\caption{Orthogonalization of B-splines to define cardinal orthogonal spline functions.}
\label{fig-spline-ortho}
\end{figure}

\begin{proof}
	The family $\{\th(\cdot-n)\}_{n \in \ZZ}$ is orthonormal if and only if
	\eq{
		\dotp{\th}{\th(\cdot-n)} = (\th \star \bar\th)(n) = \de_{0,n} \eqdef
		\choice{
			1 \qifq n=0, \\
			0 \quad\text{otherwise.}
		}
	}
	where $\bar\th(x)=\overline{\th(-x)}$ and $\de_0$ is the discrete Dirac vector.
	 
By Tonelli's theorem and Plancherel's identity, $A\in L^1(\TT)$. Its Fourier coefficients are
	\begin{align*}
		\frac{1}{2\pi}\int_0^{2\pi}A(\om)e^{-\imath n\om}\,\d\om
		&=\frac{1}{2\pi}\int_\RR |\hat\th(\om)|^2e^{-\imath n\om}\,\d\om\\
		&=\dotp{\th}{\th(\cdot+n)}.
	\end{align*}
	Thus the translates are orthonormal exactly when the Fourier coefficients of $A$ are those of the constant function 1. Uniqueness of Fourier coefficients in $L^1(\TT)$ gives $A=1$ almost everywhere.
	 
	Since $A$ is $2\pi$-periodic, normalization by the bounded function $1/\sqrt{A(\om)}$ gives $\sum_{k} |\hat \phi(\om-2\pi k)|^2 = 1$. It remains to check equality of the closed spans. Approximate the periodic multipliers $1/\sqrt A$ and $\sqrt A$ by trigonometric polynomials in $L^2(\TT)$. The bounds $a\leq A\leq b$ show that the corresponding products with $\hat\th$ and $\hat\phi$ converge in $L^2(\RR)$. Inverting the Fourier transform expresses each generator as an $L^2$ limit of finite linear combinations of translates of the other.
\end{proof}

Spline interpolation provides a typical application. For example, cubic splines are generated by translates of a box-spline function $\th$, which is a compactly supported piecewise polynomial.

                                                   
\section{Multiresolution Detail Spaces}

The detail spaces are the orthogonal complements associated with the inclusions $V_j \subset V_{j-1}$. Since these subspaces are closed,
\eq{
	\foralls j, \quad W_j \text{ is such that } V_{j-1} = V_j \oplus^\bot W_j. 
}
The following diagram shows these nested approximation spaces and their orthogonal complements:
\begin{center}
	\image{wavelets}{.8}{embeded-spaces-1d}
\end{center}
Suppose that $W_0$ admits an orthonormal basis of translates $\{\psi(\cdot-n)\}_n$. Scaling then gives
\eq{
	\{ \psi_{j,n} \}_n \text{ is a Hilbertian orthonormal basis of } W_j
	\qwhereq
	\psi_{j,n} \eqdef \frac{1}{2^{j/2}} \psi\pa{ \frac{\cdot - 2^j n}{2^j} }.
}

Orthogonal decomposition across all scales yields
\eq{
	L^2(\RR) = \bigoplus_{j=-\infty}^{j=+\infty} W_j
	=V_{j_0}\oplus^\bot\bigoplus_{j\leq j_0}W_j.
}
For every $f \in L^2(\RR)$, the corresponding expansion converges in $L^2(\RR)$:
\eq{
	f = 
	\lim_{(j_-,j_+) \rightarrow (-\infty,+\infty)} \sum_{j=j_-}^{j_+} P_{W_j} f
	= 
	\lim_{j_-\to-\infty}\left(P_{V_{j_0}}f+\sum_{j=j_-}^{j_0}P_{W_j}f\right).
}
This decomposition shows that 
\eq{
	\enscond{\psi_{j,n}}{(j,n) \in \ZZ^2}
}
is an orthonormal basis of $L^2(\RR)$, which is called a wavelet basis.
 
Stopping at a coarsest scale gives the orthonormal basis
\eq{
	\enscond{\psi_{j,n}}{j \leq j_0, n \in \ZZ} \cup 
	\enscond{\phi_{j_0,n}}{n \in \ZZ}.
}

The forward wavelet transform computes all inner products of a function $f$ with the elements of this basis.

\myfigure{
\tabdeux{
\image{wavelets}{.48}{multires-1d-haar-0}&
\image{wavelets}{.48}{multires-1d-haar-1}\\
Signal $f$ & $j=-8$ \\
\image{wavelets}{.48}{multires-1d-haar-2}&
\image{wavelets}{.48}{multires-1d-haar-3}\\
$j=-7$ & $j=-6$
}
}{ 
	1-D Haar multiresolution projection $P_{V_j}f$ of a function $f$.  
}{fig-haar-1d}



  
\paragraph{Haar wavelets.}

For the Haar multiresolution~\eqref{eq-haar-multires}, one has
\eql{\label{eq-haar-details}
	W_j = \enscond{ f }{  \foralls n \in \ZZ, \; f \text{ constant on  } [2^{j-1}n,2^{j-1}(n+1)) \qandq
	 \int_{n2^j}^{(n+1)2^j} f = 0 }.
}
One choice of mother wavelet is
\eq{
	\psi(t)=\choice{
		1 \qforq 0 \leq t < 1/2, \\
		-1 \qforq 1/2 \leq t < 1, \\
		0 \quad \text{otherwise,}
	}
}
as shown in Figure~\ref{fig-haard-scaling-wav}.


\myfigure{
\tabdeux{
\image{wavelets}{.48}{multires-1d-haar-1}& \\
$P_{V_j} f$ 